\documentclass[11pt]{article}
\usepackage{../../latex/stat4ml_style}

\title{Chapter 3: Stochastic Approximation and Optimisation \\ \large Part III --- Statistical Foundations of Deep Learning \& Reinforcement Learning}
\author{Omid Safarzadeh}
\date{}

\begin{document}
\maketitle

\begin{abstract}
\noindent
Stochastic gradient descent is not a deep-learning invention --- it is Robbins
and Monro's 1951 stochastic approximation, the statistical method for solving an
equation you can only observe through noise. Training minimises the risk, whose
gradient we never see exactly; a minibatch gives a noisy, unbiased estimate of
it, and SGD is exactly Robbins--Monro applied to the equation $\nabla R=0$. That
lineage hands us the convergence theory for free: the step sizes must shrink,
but not too fast --- $\sum a_t=\infty$ so the iterate can still travel any
distance, and $\sum a_t^2<\infty$ so the injected noise averages away. This
chapter derives those two conditions, shows the noise ball a constant step
leaves behind, and explains why the noise itself is a regulariser.
\end{abstract}

\tableofcontents

% ======================================================================
\section{Solving an equation you can only see through noise}
% ======================================================================

We want a stationary point of the risk, a $\theta$ with $\nabla R(\theta)=0$.
We cannot evaluate $\nabla R$ --- it is an expectation over the data
distribution --- but a single example or minibatch gives an \emph{unbiased}
estimate: $\mathbb{E}[\widehat{\nabla R}] = \nabla R$. This is precisely the
setting \keyterm{Robbins--Monro} solved: find the root of a function $g$ when you
observe only $g(\theta)+\text{noise}$.

\begin{definition}[Stochastic approximation]
Given noisy observations $Y_t = g(\theta_t) + \varepsilon_t$ with
$\mathbb{E}[\varepsilon_t]=0$, the Robbins--Monro iteration is
\[
  \theta_{t+1} = \theta_t - a_t\, Y_t,
\]
with a \keyterm{step-size} (or learning-rate) sequence $a_t>0$.
\end{definition}

Set $g=\nabla R$ and $Y_t=\widehat{\nabla R}(\theta_t)$ and the iteration is
stochastic gradient descent, $\theta_{t+1}=\theta_t - a_t\,\widehat{\nabla
R}(\theta_t)$. SGD is not an approximation to ``real'' gradient descent; it is
the exact statistical procedure for optimising an expectation from samples. The
only design choice is the step-size schedule, and that choice is governed by two
conditions.

% ======================================================================
\section{The two step-size conditions}
% ======================================================================

\begin{theorem}[Robbins--Monro conditions]
If the steps satisfy
\[
  \sum_{t=1}^\infty a_t = \infty \qquad\text{and}\qquad \sum_{t=1}^\infty a_t^2 < \infty,
\]
then (under mild regularity) $\theta_t$ converges to a root of $g$ almost surely.
\end{theorem}

The two conditions pull in opposite directions, and reading them is the whole
lesson. The first, $\sum a_t=\infty$, says the steps must not shrink so fast
that their total travel is bounded --- otherwise the iterate could get stuck
before reaching the root, however far away it starts. The second,
$\sum a_t^2<\infty$, says they must shrink fast enough that the accumulated
\emph{noise} --- which enters squared, through its variance --- stays finite and
averages away. The canonical schedule $a_t=1/t$ threads both: $\sum 1/t$
diverges like $\log t$, while $\sum 1/t^2 = \pi^2/6$ converges. The middle panel
of the figure shows exactly this pair of curves.

A \emph{constant} step violates the second condition, and the consequence is
visible in the left panel: the iterate does not converge to $\theta^\star$ but
settles into a \keyterm{noise ball} around it whose radius is set by the step
size. For a quadratic risk with curvature $H$ and gradient-noise variance
$\sigma^2$, the stationary variance of constant-step SGD is, to leading order,
\[
  \operatorname{Var}(\theta_\infty) \;\approx\; \frac{a\,\sigma^2}{2H},
\]
linear in the step size $a$: halve the step and you halve the noise floor. The
right panel is the same fact as rates --- $a_t=1/t$ drives the error down like
$1/t$, while a constant step plateaus at that floor.

\begin{example}[Why not just use tiny constant steps?]
Because the first condition then bites: a very small constant step has a low
noise floor but crawls, and a schedule that decays to zero is what gets both the
travel and the averaging. In practice one often holds the step constant to move
fast, then decays it to settle --- a direct reading of the two conditions.
\end{example}

% ======================================================================
\section{The noise is a feature}
% ======================================================================

There is a second reason SGD dominates in deep learning, beyond cheapness: its
noise \keyterm{regularises}. Full-batch gradient descent heads for the nearest
stationary point, sharp or flat; the stochastic gradient's fluctuations let the
iterate escape sharp, narrow minima and settle in \keyterm{flat} ones, which
empirically generalise better (a flat minimum is robust to the parameter
perturbation that a test-train shift resembles). This is \keyterm{implicit
regularisation}: no penalty is written down, yet the \emph{dynamics} bias the
solution --- toward flat minima, and, for over-parameterised linear models,
toward the minimum-norm interpolant, the same object that made the second
descent of Chapter~1 generalise.

The stability threshold sharpens the picture. On a quadratic with curvature $H$,
gradient descent with constant step $a$ contracts only if $a < 2/H$; above
$2/H$ it diverges. So the usable step size is set by the sharpest direction of
the loss, which is why learning rates and curvature (and hence normalisation,
which controls curvature) are inseparable in practice.

\begin{summarybox}
SGD is Robbins--Monro stochastic approximation applied to $\nabla R=0$: a
minibatch is an unbiased noisy gradient, and $\theta_{t+1}=\theta_t-a_t
\widehat{\nabla R}$ solves the root equation from samples. Convergence needs
$\sum a_t=\infty$ (enough travel) and $\sum a_t^2<\infty$ (noise averages away);
$a_t=1/t$ satisfies both, a constant step violates the second and leaves a noise
ball of variance $\approx a\sigma^2/2H$. The gradient noise is not only
tolerable but useful --- it biases SGD toward flat, minimum-norm solutions
(implicit regularisation) --- while the constant-step stability limit $a<2/H$
ties the usable learning rate to the loss curvature.
\end{summarybox}

\end{document}
